% Historical text retained outside the canonical input closure.
\begin{lemma}[Solved two-gap common-pair adapter]
\label{lem:appendix-two-gap-common-pair-adapter}
The reaches defined from row B in
\zcref{lem:new-two-gap-common-pair-adapter} satisfy its displayed
coordinatewise alternatives.
\end{lemma}

\begin{proof}
In the signed CE2 chart of Appendix A, the complementary far endpoints of
the two positive C traces are
\[
 p=1-(R+\alpha)=W-\alpha,
 \qquad
 q=1-(W+\delta)=R-\delta.
\]
The row predicates leave exactly
\[
 (N_+,t)\in
 \{(0,0),(0,1),(0,2),(1,0)\}.
\]
Normalize the gap edges to \(e_{5,0}\) and \(e_{0,1}\).  On the four
intervening edges choose
\[
 x_i\in(1-A_{i+1},B_i)
 \qquad(1\le i\le4).
\]
If \(N_+=0\), then \(T_1,\ldots,T_5\) are nonsupercritical immediately.
In the remaining case \((N_+,t)=(1,0)\), let \(T_s\) be supercritical.  If
\(s\ne0\), then \zcref{lem:self-midpoint} and the unique-center-midpoint
identity give
\[
 M_s\notin T_s\cup T_C.
\]
Any other open V triangle covering \(M_s\) would be adjacent and would have
positive-length trace on \(r_s\), contrary to \(d=t=0\).  Thus \(s=0\), and
\(T_1,\ldots,T_5\) are again nonsupercritical.  Gap containment and the
displayed handoffs now give
\[
 p<B_5\le1-A_5<x_4,
\]
\[
 x_4<B_4\le1-A_4<x_3,
\]
\[
 x_3<B_3\le1-A_3<x_2,
\]
\[
 x_2<B_2\le1-A_2<x_1,
\]
\[
 x_1<B_1\le1-A_1<1-q.
\]
Thus
\[
 p<x_4<x_3<x_2<x_1<1-q,
\]
and, including the two roles incident to the gap edges,
\[
 q\le A_i,
 \qquad
 p\le B_i
 \qquad(1\le i\le5).
\]

Finally, if \(T_j\) is T3-like and supports \(r_{j-1}\) or \(r_{j+1}\),
then \(t\ge1\).  The row-B identity \(N_+t=0\) gives \(N_+=0\), and hence
\[
 A_j+B_j\le1.
\]
Coordinatewise antitonicity and reflection give
\[
 C_\pm(A_j,B_j)
 \le
 \max\{C_+(p,q),C_-(p,q)\}
 \le c_{\max}(p,q)
\]
by \zcref{lem:appendix-common-pair-domination}.  This proves both assertions.
\end{proof}